\section{COMMONSENSE NORM BANK y Delphi Architecture}

El «parámetro» más importante de un modelo de juicio ético quizá no sea el hidden size, sino la dataset composition. Esta sección separa el artículo, los crowd judgments, las rules of thumb y el backbone UNICORN, y explica cómo las capacidades y los sesgos de Delphi entran al sistema a lo largo de la cadena de suministro de datos.

\subsection{Norm Bank es la confluencia de varios regímenes de datos}

La source slide enumera conjuntos de datos como Social Chemistry, ETHICS, Moral Stories y Scruples, que confluyen en COMMONSENSE NORM BANK. A la derecha se indica que hay alrededor de 1.7M people's ethical judgments on everyday situations. Esta cifra expresa annotation volume; no significa 1.7M culturas independientes ni grupos equilibrados. Además, los distintos conjuntos de datos no comparten el mismo propósito de recolección ni el mismo label schema.

\lecturefigure{slide-31-commonsense-norm-bank-sources.jpg}{Fuentes de varios conjuntos de datos de COMMONSENSE NORM BANK y 1.7M juicios}{Diapositiva V031 recuperada de la grabación oficial; Jiang et al., 2021}

\begin{knowledgebox}{Dataset Aggregation no es neutralización de valores}
Combinar datos amplía la coverage, pero no elimina por sí solo el sampling bias. Si varias fuentes se inclinan hacia el inglés, hacia usuarios de internet de Estados Unidos o hacia cierto contexto político, sumar cantidades solo estima con más precisión esa distribución. Una auditoría de alta calidad debe informar las fuentes, el periodo, los atributos demográficos, las label en conflicto y la abstention policy.
\end{knowledgebox}

\subsection{Recitar principios no equivale a saber aplicarlos}

La GPT-3 morality slide contrasta lo declarative con lo imperative. El modelo puede decir «arrojar al vecino a una licuadora está mal», pero da respuestas inestables cuando el enunciado se reformula como orden o se coloca en un contexto complejo. El callout azul enfatiza que teaching AI with principles es importante y que, al mismo tiempo, plantea problemas de dual-use y de robustness. Al leer la figura, el failure debe ubicarse en la principle application, en lugar de afirmar simplemente que el modelo «no conoce ese principio».

\lecturefigure{slide-32-gpt3-morality-declarative-imperative.jpg}{GPT-3 puede repetir principios morales, pero quizá no logre aplicarlos de forma estable}{Diapositiva V032 recuperada de la grabación oficial; Stanford CS25 Lecture 16}

\subsection{Cómo las Rules of Thumb conectan Situation y Judgment}

La diapositiva de reglas presenta primero una situation; después, el annotator escribe una rule of thumb y marca attributes como legality, social norm y moral foundation. La slide completa indica alrededor de 300,000 rules of thumb, 104k real-life situations y 12 structured attributes. Los atributos estructurados ayudan al modelo a distinguir razones distintas, como «ilegal», «descortés» o «dañino», pero también convierten la annotation ontology en el límite de sus capacidades.

\lecturefigure{slide-33-rules-of-thumb-300k.jpg}{Rules of Thumb: conexión entre situaciones reales, razones del juicio y atributos estructurales}{Diapositiva V033 recuperada de la grabación oficial; Stanford CS25 Lecture 16}

\subsection{UNICORN Backbone y la combinación de varias capacidades}

La slide de arquitectura suma el norm bank y el UNICORN commonsense model para obtener Delphi. Este signo de más no es una suma matemática, sino un pipeline de preentrenamiento / fine-tuning: UNICORN aporta una representación unificada de commonsense QA y el norm bank aporta morally salient supervision. El sistema debe manejar al mismo tiempo language understanding, commonsense inference y norm classification.

\lecturefigure{slide-34-delphi-unicorn-architecture.jpg}{COMMONSENSE NORM BANK + UNICORN forman Delphi}{Diapositiva V034 recuperada de la grabación oficial; Jiang et al., 2021}

La siguiente diapositiva responde directamente «Why Delphi is built on top of Unicorn?»: el moral reasoning requiere que language, commonsense y norms funcionen a la vez. Los ejemplos de la derecha reúnen situaciones como pedir una hamburguesa o dar consejos médicos, para mostrar que el juicio ético no es una label head independiente, sino que depende de hechos del mundo y de las consecuencias de las acciones.

\lecturefigure{slide-35-why-unicorn.jpg}{Por qué el moral reasoning requiere language, commonsense y norms}{Diapositiva V035 recuperada de la grabación oficial; Stanford CS25 Lecture 16}

\begin{warningbox}{El Backbone no puede corregir la Label Semantics}
Un encoder más potente puede mejorar la generalization dentro de la distribución, pero no determina «si el voto mayoritario es justo», «cómo se expresan culturas en conflicto» ni «cuándo debe el modelo negarse a responder». Estas cuestiones corresponden a data governance, task design y deployment policy; la representation capacity no las resuelve por sí sola.
\end{warningbox}

\subsection{Resumen de la sección}

El comportamiento de Delphi proviene de tres capas: norm data de múltiples fuentes, rules of thumb estructuradas y un commonsense backbone. Una vez que se entiende esta cadena, ni las puntuaciones altas ni los sesgos del sistema resultan misteriosos. La siguiente sección aborda la adversarial pressure posterior a la publicación abierta y examina si el desempeño dentro del dataset resiste a usuarios reales.
